\reportelement{Введение}

Цель работы --- перенести запуск приложения, PostgreSQL и Redis в Docker
Compose, добавить кэширование ответов приложения в Redis и попробовать
масштабирование сервиса.

Задачи:
\begin{itemize}
  \item вынести персональные параметры в файл \cmd|.env| и описать стек
        в \cmd|docker-compose.yml|;
  \item настроить порядок запуска сервисов через \cmd|healthcheck| и
        \cmd|depends_on|;
  \item добавить сервис Redis и кэширование эндпоинта
        \cmd|/weather/{city}|;
  \item проверить работу кэша, сетей и томов;
  \item запустить несколько экземпляров приложения и убедиться, что они
        используют общий кэш.
\end{itemize}

Работа выполнялась на сервере linux-lab, Docker версии 29.1.3 и Docker
Compose версии 2.40.3~\cite{docker-docs, compose-docs}.

\reportelement{Персональные параметры}

Персональные параметры приведены в таблице~\ref{tab:parameters}.

\begin{table}[H]
  \caption{Персональные параметры}\label{tab:parameters}
  \centering
  \begin{tabularx}{\textwidth}{|p{0.52\textwidth}|X|}
    \hline
    Параметр & Значение\\ \hline
    Номер ИСУ & 467697\\ \hline
    Персональный порт (PostgreSQL, \cmd|PG_PORT|) & 5527\\ \hline
    Персональный порт (Redis, \cmd|REDIS_PORT|) & 6467\\ \hline
    Персональный порт (приложение, \cmd|APP_PORT|) & 8097\\ \hline
    Персональный префикс & ter-stepanyan\\ \hline
    Секретный ключ & secret97ter-stepanyan\\ \hline
    Имя базы данных & ter-stepanyan97\\ \hline
  \end{tabularx}
\end{table}


Репозиторий с результатом (ветка \cmd|lab-4|): \url{\RepositoryURL}.

\chapter{Основная часть}

\section{Этап 1. Подготовка проекта для Docker Compose}

От ветки \cmd|lab-3| создана ветка \cmd|lab-4|. В корне проекта
\cmd|solutions/assist-py| создан файл \cmd|.env| с персональными
переменными (рисунок~\ref{fig:1-1}).  Файл содержит пароли, поэтому он добавлен в
\cmd|.gitignore|, а в репозиторий помещён \cmd|.env.example| с теми
же переменными без реальных паролей (приложение~Б).

\reportscreenshot{fig-1-1.png}{Заполнение файла .env}{fig:1-1}

Создан файл \cmd|docker-compose.yml| с двумя сервисами: \cmd|db|
(PostgreSQL) и \cmd|app| (приложение). Полное содержимое итогового
файла приведено в приложении~А. Основные решения:
\begin{itemize}
  \item PostgreSQL запускается командой \cmd|postgres -p
        ${PG_PORT}|, поэтому внутри контейнера он слушает
        персональный порт 5527;
  \item данные базы хранятся в именованном томе
        \cmd|ter-stepanyan-pgdata|;
  \item у \cmd|db| задан \cmd|healthcheck| с командой
        \cmd|pg_isready|, которая проверяет, что база готова
        принимать подключения на порту 5527;
  \item приложение собирается из \cmd|Dockerfile.optimized|, слушает
        порт 8097, и этот порт пробрасывается на хост;
  \item приложение подключается к базе по строке
        \cmd|DATABASE_URL|, в которой хостом указано имя сервиса
        \cmd|db|, а не \cmd|localhost|: у каждого контейнера свой
        сетевой стек, и \cmd|localhost| внутри контейнера указывает на
        сам контейнер;
  \item папка \cmd|config| подключена как bind mount в режиме только
        для чтения, как в ЛР~№3;
  \item приложение запускается только после того, как база пройдёт
        проверку здоровья (\cmd|depends_on| с
        \cmd|condition: service_healthy|).
\end{itemize}

Запуск всех сервисов командой \cmd|docker compose up -d --build|
показан на рисунке~\ref{fig:1-2}. При первом запуске Compose добавил к
именам сети и тома префикс проекта \cmd|assist-py_| (имя папки
проекта). Чтобы сеть и том назывались ровно
\cmd|ter-stepanyan-backend| и \cmd|ter-stepanyan-pgdata|, в
описание сети и тома позже добавлен ключ \cmd|name|.

\reportscreenshot{fig-1-2.png}{Запуск сервисов через Docker Compose}{fig:1-2}

Статус сервисов и логи приложения представлены на
рисунке~\ref{fig:1-3}. Оба контейнера в состоянии \cmd|healthy|,
приложение доступно на персональном порту 8097.

\reportscreenshot{fig-1-3.png}{Статус сервисов и логи приложения}{fig:1-3}

В логах базы данных (рисунок~\ref{fig:1-4}) видно, что PostgreSQL
слушает персональный порт 5527.

\reportscreenshot{fig-1-4.png}{Логи базы данных}{fig:1-4}

Приложение отвечает на запросы \cmd|/| и \cmd|/health|
(рисунок~\ref{fig:1-5}), в ответе указан владелец
\cmd|ter-stepanyan|.

\reportscreenshot{fig-1-5.png}{Проверка ответа приложения}{fig:1-5}

\section{Этап 2. Добавление Redis для кэширования погоды}

В \cmd|docker-compose.yml| добавлен сервис \cmd|redis|
(рисунок~\ref{fig:2-1}) из образа \cmd|redis:7-alpine|. Он подключён к
той же сети \cmd|ter-stepanyan-backend|, а его порт опубликован на хосте
под персональным номером 6467. Проверка здоровья выполняется командой
\cmd|redis-cli ping|.

\reportscreenshot{fig-2-1.png}{Сервис redis в docker-compose.yml}{fig:2-1}

Приложение теперь зависит и от базы, и от Redis
(рисунок~\ref{fig:2-2}): Compose запустит его только после того, как оба
сервиса станут \cmd|healthy|.

\reportscreenshot{fig-2-2.png}{Зависимость приложения от db и redis}{fig:2-2}

В \cmd|app/main.py| добавлено подключение к Redis
(рисунок~\ref{fig:2-3}), а в \cmd|requirements.txt| --- библиотека
\cmd|redis|. Адрес Redis и время жизни кэша приложение берёт из
переменных окружения \cmd|REDIS_URL| и \cmd|CACHE_TTL|.

\reportscreenshot{fig-2-3.png}{Подключение к Redis в main.py}{fig:2-3}

Эндпоинт \cmd|/weather/{city}| работает так:
\begin{enumerate}
  \item приложение формирует ключ вида
        \cmd|weather:ter-stepanyan:moscow| и запрашивает его в Redis;
  \item если значение найдено, ответ возвращается из кэша с полем
        \cmd|"source": "cache"| без обращения к внешнему API;
  \item если значения нет, приложение запрашивает погоду у API
        open-meteo, возвращает ответ с \cmd|"source": "api"| и
        сохраняет его в Redis командой \cmd|SETEX| на
        \cmd|CACHE_TTL| = 120 секунд.
\end{enumerate}

Так повторные запросы по одному городу не нагружают внешний API и
обрабатываются быстрее, а через 120 секунд Redis сам удаляет ключ, и
данные обновляются.

Также добавлен эндпоинт \cmd|/cache/stats| (рисунок~\ref{fig:2-4}). Он
находит в Redis все ключи погоды текущего владельца и возвращает их
количество и список городов.

\reportscreenshot{fig-2-4.png}{Эндпоинт статистики кэша}{fig:2-4}

В переменные окружения приложения добавлены \cmd|REDIS_URL| и
\cmd|CACHE_TTL| (рисунок~\ref{fig:2-5}). В \cmd|REDIS_URL| хостом
указано имя сервиса \cmd|redis|, а портом --- порт, на котором Redis
слушает внутри своего контейнера. Проброс порта 6467 на хост для связи
приложения с Redis не нужен: контейнеры общаются напрямую по
пользовательской сети, а встроенный DNS Docker сопоставляет имя сервиса
с IP-адресом контейнера. Проброшенный порт нужен только для доступа к
Redis с хоста.

\reportscreenshot{fig-2-5.png}{Переменные окружения приложения}{fig:2-5}

\section{Этап 3. Запуск и проверка полного стека}

После пересборки все три сервиса запущены и находятся в состоянии
\cmd|healthy| (рисунок~\ref{fig:3-1}): \cmd|ter-stepanyan-postgres|,
\cmd|ter-stepanyan-assistant| и \cmd|ter-stepanyan-redis|. По
выводу Compose видно, что \cmd|postgres| и \cmd|redis| стали
\cmd|Healthy| через 5,6~с, а приложение запущено после них, через
5,8~с.

\reportscreenshot{fig-3-1.png}{Запуск полного стека}{fig:3-1}

Логи всех сервисов представлены на рисунке~\ref{fig:3-2}. PostgreSQL и
Redis сообщают о готовности принимать подключения
(\cmd|database system is ready to accept connections| и
\cmd|Ready to accept connections|), после чего стартует приложение.
Предупреждение Redis \cmd|Memory overcommit must be enabled| относится
к настройке ядра хоста и не влияет на работу кэша.

\reportscreenshot{fig-3-2.png}{Логи сервисов}{fig:3-2}

Первые строки лога приложения (рисунок~\ref{fig:3-3}) показывают запуск
uvicorn на порту 8097. Строки \cmd|Waiting for application startup|
относятся к подготовке самого приложения FastAPI. Ожидание базы и Redis
выполняет Compose ещё до создания процесса приложения, поэтому в логах
приложения оно не отображается.

\reportscreenshot{fig-3-3.png}{Начало лога приложения}{fig:3-3}

Проверка кэширования показана на рисунке~\ref{fig:3-4}. Первый запрос
\cmd|/weather/Moscow| вернул данные с \cmd|"source": "api"|, второй
такой же запрос --- с \cmd|"source": "cache"|. Эндпоинт
\cmd|/cache/stats| показал владельца \cmd|ter-stepanyan| и город
\cmd|moscow| в кэше. Команда \cmd|redis-cli KEYS| нашла ключ
\cmd|weather:ter-stepanyan:moscow|, а команда \cmd|TTL| показала
оставшееся время жизни ключа: сначала 11, затем 5 секунд. Ответ
\cmd|-2| означает, что ключ уже удалён, так как истекло его время
жизни (ответ \cmd|-1| означал бы ключ без срока жизни).

\reportscreenshot{fig-3-4.png}{Проверка кэширования}{fig:3-4}

Список сетей с моим префиксом приведён на рисунке~\ref{fig:3-5}.

\reportscreenshot{fig-3-5.png}{Список сетей}{fig:3-5}

Команда \cmd|docker network inspect ter-stepanyan-backend|
(рисунок~\ref{fig:3-6}) показывает, что все три контейнера подключены к
одной сети и получили адреса из её подсети \cmd|172.18.0.0/16|.

\reportscreenshot{fig-3-6.png}{Три контейнера в сети ter-stepanyan-backend}{fig:3-6}

Список томов приведён на рисунке~\ref{fig:3-7}: данные PostgreSQL
хранятся в именованном томе \cmd|ter-stepanyan-pgdata|.

\reportscreenshot{fig-3-7.png}{Список томов}{fig:3-7}

\section{Этап 4. Масштабирование сервиса}

Для масштабирования создан файл \cmd|docker-compose.scale.yml|
(приложение~В). У сервиса \cmd|app| убран \cmd|container_name|,
так как имя контейнера должно быть уникальным, а в \cmd|ports|
указан только порт внутри контейнера \cmd|${APP_PORT}|. Порт на
хосте Docker назначает каждому экземпляру сам, поэтому три копии не
конфликтуют за один порт хоста. Внутри все экземпляры слушают порт 8097,
что возможно, потому что у каждого контейнера свой сетевой стек.

PostgreSQL и Redis в этой конфигурации описаны по шаблону задания:
порты на хост не публикуются, и доступ к ним есть только у контейнеров
внутри сети. Снаружи доступны только экземпляры приложения.

Запуск трёх экземпляров командой
\cmd|docker compose -f docker-compose.scale.yml up -d --build --scale app=3|
показан на рисунке~\ref{fig:4-1}. Экземпляры
\cmd|assist-py-app-1|, \cmd|-2| и \cmd|-3| получили порты
хоста 32777, 32779 и 32778.

\reportscreenshot{fig-4-1.png}{Три экземпляра приложения}{fig:4-1}

Все три экземпляра отвечают на \cmd|/health| по своим портам
(рисунок~\ref{fig:4-2}).

\reportscreenshot{fig-4-2.png}{Проверка здоровья экземпляров}{fig:4-2}

Проверка общего кэша показана на рисунке~\ref{fig:4-3}. Запрос
\cmd|/weather/London| к первому экземпляру ушёл во внешний API
(\cmd|"source": "api"|), а такой же запрос ко второму экземпляру
вернул данные из кэша (\cmd|"source": "cache"|), хотя это другой
контейнер. Третий экземпляр в \cmd|/cache/stats| показывает
\cmd|london| в списке закэшированных городов. Это работает, потому что
приложение не хранит состояние внутри себя: заметки находятся в
PostgreSQL, а кэш --- в Redis, общих для всех экземпляров.

\reportscreenshot{fig-4-3.png}{Общий кэш Redis для всех экземпляров}{fig:4-3}

После проверки масштабированная конфигурация остановлена
(рисунок~\ref{fig:4-4}), и запущена основная (рисунок~\ref{fig:4-5}).
Compose снова дождался состояния \cmd|Healthy| у \cmd|redis| и
\cmd|postgres| и только после этого запустил приложение.

\reportscreenshot{fig-4-4.png}{Остановка масштабированной конфигурации}{fig:4-4}

\reportscreenshot{fig-4-5.png}{Запуск основной конфигурации}{fig:4-5}

\section{Этап 5. Финальный артефакт}

\reportscreenshot{fig-5-1.png}{Финальный артефакт}{fig:5-1}

\reportelement{Заключение}

В ходе работы запуск приложения, PostgreSQL и Redis перенесён в Docker
Compose: весь стек описан в одном файле и поднимается одной командой, а
персональные параметры вынесены в \cmd|.env|. С помощью
\cmd|healthcheck| и \cmd|depends_on| настроен порядок запуска:
приложение стартует только после того, как база и Redis готовы принимать
подключения.

В приложение добавлено кэширование погоды в Redis с временем жизни
120~секунд: повторный запрос по тому же городу возвращается из кэша без
обращения к внешнему API. Проверено, что контейнеры общаются по
пользовательской сети по именам сервисов, а проброс портов нужен только
для доступа с хоста.

При масштабировании три экземпляра приложения работали одновременно и
использовали общий кэш Redis. Это показывает, что приложение можно
масштабировать горизонтально, потому что его состояние вынесено в
отдельные сервисы.



\startreportappendices

\reportappendix{Содержимое docker-compose.yml}

\begin{code}
services:
  db:
    image: postgres:15
    container_name: ter-stepanyan-postgres
    command: postgres -p ${PG_PORT}
    environment:
      POSTGRES_USER: ${POSTGRES_USER}
      POSTGRES_PASSWORD: ${POSTGRES_PASSWORD}
      POSTGRES_DB: ${POSTGRES_DB}
    volumes:
      - ter-stepanyan-pgdata:/var/lib/postgresql/data
    networks:
      - ter-stepanyan-backend
    healthcheck:
      test: ["CMD-SHELL", "pg_isready -U ${POSTGRES_USER} -d ${POSTGRES_DB} -p ${PG_PORT}"]
      interval: 30s
      timeout: 5s
      start_period: 10s
      retries: 3

  app:
    build:
      context: .
      dockerfile: Dockerfile.optimized
    container_name: ter-stepanyan-assistant
    ports:
      - "${APP_PORT}:${APP_PORT}"
    environment:
      APP_OWNER: ${APP_OWNER}
      APP_SECRET: ${APP_SECRET}
      DATABASE_URL: postgresql://${POSTGRES_USER}:${POSTGRES_PASSWORD}@db:${PG_PORT}/${POSTGRES_DB}
      REDIS_URL: redis://redis:6379/0
      CACHE_TTL: ${CACHE_TTL}
    volumes:
      - ./config:/app/config:ro
    depends_on:
      db:
        condition: service_healthy
      redis:
        condition: service_healthy
    networks:
      - ter-stepanyan-backend

  redis:
    image: redis:7-alpine
    container_name: ter-stepanyan-redis
    ports:
      - "${REDIS_PORT}:6379"
    networks:
      - ter-stepanyan-backend
    healthcheck:
      test: ["CMD", "redis-cli", "ping"]
      interval: 5s
      timeout: 3s
      retries: 5

networks:
  ter-stepanyan-backend:
    name: ter-stepanyan-backend
    driver: bridge

volumes:
  ter-stepanyan-pgdata:
    name: ter-stepanyan-pgdata
\end{code}

\reportappendix{Содержимое .env.example}

\begin{code}
APP_OWNER=ter-stepanyan
APP_SECRET=
APP_PORT=8097
POSTGRES_USER=user_97
POSTGRES_PASSWORD=
POSTGRES_DB=assistant_97
PG_PORT=5527
REDIS_PORT=6467
CACHE_TTL=120
\end{code}

\reportappendix{Содержимое docker-compose.scale.yml}

\begin{code}
services:
  db:
    image: postgres:15
    container_name: ter-stepanyan-postgres
    environment:
      POSTGRES_USER: ${POSTGRES_USER}
      POSTGRES_PASSWORD: ${POSTGRES_PASSWORD}
      POSTGRES_DB: ${POSTGRES_DB}
    volumes:
      - ter-stepanyan-pgdata:/var/lib/postgresql/data
    networks:
      - ter-stepanyan-backend
    healthcheck:
      test: ["CMD-SHELL", "pg_isready -U ${POSTGRES_USER}"]
      interval: 5s
      timeout: 3s
      retries: 5

  redis:
    image: redis:7-alpine
    container_name: ter-stepanyan-redis
    networks:
      - ter-stepanyan-backend
    healthcheck:
      test: ["CMD", "redis-cli", "ping"]
      interval: 5s
      timeout: 3s
      retries: 5

  app:
    build:
      context: .
      dockerfile: Dockerfile.optimized
    ports:
      - "${APP_PORT}"
    environment:
      APP_OWNER: ${APP_OWNER}
      APP_SECRET: ${APP_SECRET}
      DATABASE_URL: postgresql://${POSTGRES_USER}:${POSTGRES_PASSWORD}@db:5432/${POSTGRES_DB}
      REDIS_URL: redis://redis:6379/0
      CACHE_TTL: ${CACHE_TTL}
    depends_on:
      db:
        condition: service_healthy
      redis:
        condition: service_healthy
    networks:
      - ter-stepanyan-backend

volumes:
  ter-stepanyan-pgdata:

networks:
  ter-stepanyan-backend:
    driver: bridge
\end{code}
